\chapter{Eye Tracker}
\label{sec:eye tracker}

\section{Grundlegende Hardwarevoraussetzungen}

<<<<<<< HEAD
In diesem Kapitel wird ein grundlegender Einblick in die Hardwarevoraussetzungen für ein Eye Tracking System bereitgestellt. Daraufhin wird eine Marktanalyse durchgeführt, um verfügbare Eye Tracker auf deren Anwendbarkeit, Integration in CARLA und Kosten zu untersuchen.

Eye Tracking Systeme bestimmen die Mitte der Pupille, leiten die Rotation des Augapfels ab und bestimmen dann mithilfe von mathematischen Algorithmen die Blickposition \cite{lit_website_eyeware_eyetracker_blog}. Dabei wird zwischen remote Systemen und Head Mounted Systemen unterschieden. Bei den remote Systemen sitzt das Eye Tracking System nicht auf dem Kopf der Testperson. Es ist an einer von der Testperson entfernten Stelle positioniert und erfasst von dort aus die Blickposition. Häufig werden solche remote Systeme verwendet, wenn die Blickposition der Testperson auf einen Bildschirm untersucht werden soll. Dadurch ist die Erfassung meist auf den jeweiligen Bildschirm beschränkt. Die Erfassung von Objekten in der realen dreidimensionalen Welt, welche über den Bildschirm hinausgeht, ist in der Regel nicht möglich. Es gibt allerdings ein paar Anbieter, die remote Systeme anbieten, welche auch Blickpositionen in der realen dreidimensionalen Umgebung erfassen. Bei den  tragbaren Eye Tracking Systemen hingegen wird der Testperson das System aufgesetzt. Solche Systeme erfassen die Augenbewegung dadurch aus nächster Nähe. Ein Head Tracking System eignet sich insbesondere für Tracking-Experimente in der realen Welt. \cite{lit_website_eyeware_eyetracker_blog},\cite{lit_website_bitbrain}

Die minimalen Hardwarevoraussetzungen für einen Eye Tracker belaufen sich auf eine Kamera, eine Lichtquelle und einen Computer. In teureren Systemen werden teilweise zwei Kameras eingesetzt. Dadurch kann sich die Tracking-Qualität verbessern, diese Systeme sind meist aber auch komplexer \cite{lit_website_eyeware_eyetracker_blog}. Zudem wird hier auch Software benötigt, welche mehrere Kameras gleichzeitig auswerten kann.
In diesem Projekt soll der Blick auf eine Leinwand erfasst werden. Aus diesem Grund reicht ein remote System aus. Es wird eine einzelne Kamera verwendet, da die Auswahl der in Frage kommenden Softwarelösungen sonst zu sehr eingeschränkt wird. 
Unabhängig vom eingesetzten System gilt, dass die Genauigkeit gesteigert werden kann, indem eine hochauflösende Kamera genutzt wird \cite{lit_website_bitbrain}. Ergebnisse hängen allerdings auch stark mit dem Abstand der Testperson zur Kamera ab. Sitzt die Testperson in geringem Abstand zu der Kamera, dann ist auch eine kleinere Auflösung ausreichend. In unseren Experimenten wurde eine Kamera mit einer Auflösung von 1920x1080 Pixel eingesetzt. Diese Auflösung hat sich in Kombination mit einem Abstand von unter einem Meter zur Kamera als ausreichend erwiesen, sodass der Eye Tracker die Augenbewegung erfassen konnte. Auch Kameras mit höherer Bildrate können bessere Ergebnisse liefern. Allerdings reicht eine \Gls{Bildwiederholrate} von 30 Bildern pro Sekunde für die Erfassung von Augenbewegungen in vielen Anwendungen aus \cite{lit_tracking_framerate}. Kameras haben im Vergleich zum menschlichen Auge die Fähigkeit, infrarotes ({\acs{IR}}) Licht zu erfassen. Digitalkameras und Handykameras können je nach eingesetzter Sensortechnologie bereits {\acs{IR}}-Strahlung wahrnehmen. Oft sind allerdings Filter eingebaut, die {\acs{IR}}-Strahlung größtenteils herausfiltern, da die Kameras die menschliche Sicht abbilden sollen \cite{lit_website_eyeware_eyetracker_blog}. Es gibt allerdings spezielle {\acs{IR}}-Kameras, deren Sensor auf einer anderen Technologie basiert, sodass diese besonders im {\acs{IR}}-Bereich sensitiv sind. Einsatzgebiete sind bspw. Wärmebildkameras oder Nachtsichtkameras. Im Kontext von Eye Tracking Systemen können solche Kameras auch bei schlechten Lichtverhältnissen, im Vergleich zu einer Kamera, welche nur im menschlich sichtbaren Frequenzbereich operiert, noch verwertbare Bildinformationen erhalten \cite{lit_article_infrared}.
Durch das Anbringen von Infrarotlichtquellen und dem Einsatz einer solchen Infrarotkamera können auch bei schlechteren Lichtverhältnissen verlässliche Ergebnisse erzielt werden \cite{lit_article_stateoftheart_eyetracker}.

Auf die genaue Funktionsweise von Eye Tracking Systemen und den dafür eingesetzten mathematischen Algorithmen soll in dieser Projektarbeit nicht eingegangen werden, da es nicht die Zielsetzung ist, ein solches System selbst zu entwickeln. Für die vorliegende Arbeit ist lediglich relevant, dass die Ergebnisqualität des Eye Tracking System stark von der Art des eingesetzten Systems, den Lichtverhältnissen und dem Abstand der Testperson zur Kamera abhängig ist.

\section{Marktanalyse} 

Wie oben bereits beschrieben wird eine Marktanalyse von bereits verfügbaren Eye Trackern durchgeführt. Die Kriterien, die das System erfüllen soll, wurden bereits in vorherigen Abschnitten erläutert. 

Ein erster Softwarekandidat ist ein \Gls{Open Source} Tracking Projekt auf GitHub \cite{lit_code_github_facegazetracker}. Die Ausgabe dieses Trackers ist in Abbildung \ref{fig:old_tracker_face_gaze_webcam} zu erkennen. Das System orientiert sich hauptsächlich an den Augen und der Nase der Testperson. Dadurch kann es bestimmen, in welche Richtung der Kopf einer betrachteten Person geneigt ist. Zudem kann die Software das Blinzeln der Augen erkennen.

=======
In diesem Kapitel wird ein grundlegender Einblick in Hardwarevoraussetzungen für ein Eye Tracking System bereitgestellt. Daraufhin wird eine Marktanalyse durchgeführt, um verfügbare Eye Tracker auf deren Anwendbarkeit, Integration in CARLA und Kosten zu untersuchen.
Eye Tracking Systeme bestimmen die Mitte der Pupille, leiten die Rotation des Augapfels ab und bestimmen dann mithilfe von mathematischen Algorithmen die Blickposition \cite{lit_website_eyeware_eyetracker_blog}. Dabei wird zwischen remote Systemen und Head Mounted Systemen unterschieden. Bei den remote Systemen sitzt das Eye Tracking System nicht auf dem Kopf der Testperson. Es ist an einer von der Testperson entfernten Stelle positioniert und erfasst von dort aus die Blickposition. Häufig werden solche remote Systemen verwendet, wenn die Blickposition der Testperson auf einen Bildschirm untersucht werden soll. Dadurch ist die Erfassung meist auf den jeweiligen Bildschirm beschränkt. Die Erfassung von Objekten in der realen dreidimensionalen Welt, welche über den Bildschirm hinaus geht, ist in der Regel nicht möglich. Es ist aber nicht unmöglich, es gibt ein paar Anbieter die remote Systeme anbieten, welche auch Blickpositionen in der realen dreidimensionalen Umgebung erfassen. Bei den  tragbaren Eye Tracking Systemen hingegen wird der Testperson das System aufgesetzt. Solche Systeme erfassen die Augenbewegung dadurch aus nächster Nähe. Ein Head Tracking System ist für Tracking Experimente in der realen Welt besonders gut geeignet. \cite{lit_website_eyeware_eyetracker_blog},\cite{lit_website_bitbrain}

Die minimalen Hardwarevoraussetzungen für einen Eye Tracker belaufen sich auf eine Kamera, eine Lichtquelle und einen Computer. Als Kamera genügt eine handelsübliche Webcam. In teureren Systemen werden teilweise zwei Kameras eingesetzt. Dadurch kann sich die Tracking Qualität verbessern, diese Systeme sind meist aber auch komplexer \cite{lit_website_eyeware_eyetracker_blog}. Zudem wird hier auch Software benötigt, welche mehrere Kameras gleichzeitig auswerten kann.
In diesem Projekt soll der Blick auf eine Leinwand erfasst werden. Aus diesem Grund reicht ein remote System aus. Es wird eine einzelne Kamera verwendet, da die Auswahl der in Frage kommenden Softwarelösungen sonst zu sehr eingeschränkt wird. 
Unabhängig vom eingesetzt System gilt, dass die Genauigkeit gesteigert werden kann, indem eine hochauflösende Kamera genutzt wird. Ergebnisse hängen allerdings auch stark mit dem Abstand der Testperson zur Kamera ab. Sitzt die Testperson in geringem Abstand zu der Kamera, dann ist auch eine kleinere Auflösung ausreichend. In unseren Experimenten wurde eine Kamera mit einer Auflösung von 1920x1080 Pixel eingesetzt. Diese Auflösung hat sich in Kombination mit einem Abstand von unter einem Meter zur Kamera als ausreichend erwiesen. Auch Kameras mit höherer Bildrate können bessere Ergebnisse liefern. Allerdings reicht eine \Gls{Bildwiederholrate} von 30 Bildern pro Sekunde für die Erfassung von Augenbewegungen in vielen Anwendungen aus \cite{lit_tracking_framerate}. Kameras haben im Vergleich zum menschlichen Auge die Fähigkeit infrarotes ({\acs{IR}}) Licht zu erfassen. Digitalkameras und Handykameras können je nach eingesetzter Sensortechnologie bereits {\acs{IR}}-Strahlung wahrnehmen. Oft sind allerdings Filter eingebaut, die {\acs{IR}}-Strahlung größtenteils herausfiltern, da die Kameras die menschliche Sicht abbilden sollen \cite{lit_website_eyeware_eyetracker_blog}. Es gibt allerdings spezielle {\acs{IR}}-Kameras, deren Sensor auf einer anderen Technologie basiert, sodass diese besonders im {\acs{IR}}-Bereich sensitiv sind. Einsatzgebiete sind bspw. Wärmebildkameras oder Nachtsichtkameras. Im Kontext von Eye Tracking Systemen können solche Kameras selbst bei schlechten Lichtverhältnissen , im Vergleich zu einer Kamera, welche nur im menschlich sichtbaren Frequenzbereich operiert, noch verwertbare Bildinformation erhalten \cite{lit_article_infrared}.
Durch das Anbringen von Infrarotlampen und dem Einsatz einer solchen Infrarotkamera können auch bei schlechteren Lichtverhältnissen verlässliche Ergebnisse erzielt werden \cite{lit_article_stateoftheart_eyetracker}.

Auf die genaue Funktionsweise von Eye Tracking Systemen und den dafür eingesetzten mathematischen Algorithmen soll in dieser Projektarbeit nicht eingegangen werden, da es nicht die Zielsetzung ist, ein solches System selbst zu entwickeln. Für die vorliegende Arbeit ist lediglich relevant, dass die Ergebnisqualität des Eye Tracking System stark von der Art des eingesetzten Systems, den Lichtverhältnissen und vom Abstand der Testperson zur Kamera abhängig ist.

\section{Marktanalyse} 

Da die Entwicklung eines eigenen Eye Tracking Algorithmus den zeitlichen Rahmen der Projektarbeit übersteigt, wird eine Marktanalyse von bereits verfügbaren Eye Trackern durchgeführt. Die Kriterien, die das System erfüllen soll, wurden bereits in vorherigen Abschnitten erläutert. 

Ein erster Softwarekandidat ist ein \Gls{Open Source} Tracking Projekt auf GitHub \cite{lit_code_github_facegazetracker}. Die Ausgabe dieses Trackers ist in Abbildung \ref{fig:old_tracker_face_gaze_webcam} zu erkennen. Das System orientiert sich hauptsächlich an den Augen und der Nase der Testperson. Dadurch kann es bestimmen, in welche Richtung der Kopf einer betrachteten Person geneigt ist. Zudem kann die Software das Blinzeln der Augen erkennen.

Die Software schätzt nicht die Blickposition der Testperson auf dem Bildschirm, sondern gibt nur an, in welcher Lage sich der Kopf befindet. Diese Daten könnten aber als Grundlage für eine Weiterentwicklung dienen. Da allerdings nur die Ausrichtung des Gesichts und nicht der Pupillen in die Ausgabe einfließt, werden keine genauen Ergebnisse durch diese Herangehensweise erwartet.
>>>>>>> fd0befd8e95e2bae230d21ece8a1862eada6b577

\begin{center}
	\includegraphics[scale=0.4]{old_tracker_face_gaze_webcam}
	\captionof{figure}{Ausgabe des Face-Gaze-Trackers, Quelle: \cite{lit_code_github_facegazetracker}} \label{fig:old_tracker_face_gaze_webcam}
\end{center}

<<<<<<< HEAD
=======
Eine weitere Option stellt das open source Projekt "Eyegaze Tracker" dar. Diese Software erfasst nicht die Gesichtsausrichtung, sondern die Pupillen. Dadurch kann sie Aussagen über die grobe Blickrichtung treffen. So wie in Abbildung \ref{fig:old_tracker_eyegaze_webcam} zu erkennen ist, kann bspw. die Aussage getroffen werden, dass die Testperson nach links schaut. Auch hier wird keine Blickposition auf dem Bildschirm geschätzt. Allerdings werden die Pupillen erfasst und die grobe Blickrichtung bestimmt. Die Blickposition kann gegebenenfalls mit den vorhandenen Daten durch eine Weiterentwicklung des vorhandenen open source Codes bestimmt werden. Dies würde allerdings einen erheblichen Entwicklungsaufwand bedeuten. 
>>>>>>> fd0befd8e95e2bae230d21ece8a1862eada6b577

Die Software schätzt nicht die Blickposition der Testperson auf dem Bildschirm, sondern gibt nur an, in welcher Lage sich der Kopf befindet. Diese Daten könnten aber als Grundlage für eine Weiterentwicklung dienen. Da allerdings nur die Ausrichtung des Gesichts und nicht der Pupillen in die Ausgabe einfließt, werden keine genauen Ergebnisse durch diese Herangehensweise erwartet.

<<<<<<< HEAD

Eine weitere Option stellt das open source Projekt 'Eyegaze Tracker' dar. Diese Software erfasst nicht die Gesichtsausrichtung, sondern die Pupillen. Dadurch kann sie eine qualitative Aussage über die Blickrichtung treffen. So wie in Abbildung \ref{fig:old_tracker_eyegaze_webcam} zu erkennen ist, kann bspw. die Aussage getroffen werden, dass die Testperson nach links schaut. Auch hier wird keine Blickposition auf dem Bildschirm geschätzt. Allerdings werden die Pupillen erfasst und die grobe Blickrichtung bestimmt. Die Blickposition kann gegebenenfalls mit den vorhandenen Daten durch eine Weiterentwicklung des vorhandenen open source Codes bestimmt werden. Dies würde allerdings einen erheblichen Entwicklungsaufwand bedeuten. 


\begin{center}
	\includegraphics[scale=0.4]{old_tracker_eyegaze_webcam}
	\captionof{figure}{Ausgabe des Eyegaze-Trackers, Quelle: \cite{lit_code_github_gazetracking}} \label{fig:old_tracker_eyegaze_webcam}
\end{center}


Neben den oben genannten Projekten wurden auch noch einige weitere auf ihre Einsetzbarkeit in diesem Projekt untersucht. Diese hatten größtenteils ähnliche Eigenschaften und Einschränkungen. Eye Tracker die einen größeren Funktionsumfang bieten, wie beispielsweise das System 'Tobii Eye Tracker' haben den Nachteil, dass sie einen, je nach Budget, hohen finanziellen Aufwand bedeuten. Es wurde sich auf Grundlage umfangreicher Recherche und experimenteller Erprobung für den Beam Eye Tracker von Eyeware Tech entschieden. Dieser hat die in dieser Arbeit definierten Anforderungen erfüllt und soll als Basis für die weitere Entwicklung dienen.

\newpage

\section{Beam Eye Tracker}

Der Beam Eye Tracker ist für Gaming-Anwendungen und dem Einsatz von einfacher Hardware wie normalen Webcams entwickelt. Neben der Schätzung der Blickposition auf einen Punkt auf dem Bildschirm verfügt das Programm über eine Kopfpositionserkennung. Durch die Funktion 'Gaming Erweiterungen' teilt der Beam Eye Tracker die Trackingdaten über seine {\ac{API}} mit anderen Programmen.

Die Beam {\ac{API}} kann über die Programmiersprachen C++, C, C\# und Python angesprochen werden \cite{lit_code_beam_sdk}. Im Rahmen dieser Projektarbeit wird die {\ac{API}} mit Pyhton angesprochen. Verschiedene Klassen in der {\ac{API}} ermöglichen eine individuelle Einbindung der Eye Tracking Daten in eine selbst entwickelte Anwendung. Relevant für diese Projektarbeit ist die Klasse \glqq{}screen\_gaze\_info\grqq{}. Diese beinhaltet die x- und y- Koordinaten und die Angaben, ob die Augen erkannt werden und wie gut die Schätzung ist.

Durch die anwenderfreundliche {\ac{API}} des Beam {\ac{SDK}} und den niedrigen Preis von 29,99€ des Eye Trackers, fiel die Entscheidung, diesen für das Projekt zu nutzen. In Rücksprache mit den Beam-Entwicklern wurde der Beam Eye Tracker für dieses Projekt unentgeltlich zur Verfügung gestellt.
=======
Neben den oben genannten Projekten wurden auch noch einige weitere auf ihre Einsetzbarkeit in diesem Projekt untersucht. Diese hatten größtenteils ähnliche Eigenschaften und Einschränkungen. Eye Tracker die einen größeren Funktionsumfang bieten wie beispielsweise das System "Tobii Eye Tracker" haben den Nachteil, dass sie einen,je nach Budget, hohen finanziellen Aufwand bedeuten. Es wurde sich auf Grundlage umfangreicher Recherche und experimenteller Erprobung für den Beam Eye Tracker von Eyeware Tech entschieden. Dieser hat die in dieser Arbeit definierten Anforderungen erfüllt und soll als Basis für die weitere Entwicklung dienen.

\section{Beam Eye Tracker}

Der Beam Eye Tracker ist für Gaming-Anwendungen und dem Einsatz von einfacher Hardware wie normalen Webcams entwickelt. Neben der Schätzung der Blickposition auf einen Punkt auf dem Bildschirm verfügt das Programm über eine Kopfpositionserkennung. Durch die Funktion "Gaming Erweiterungen" teilt der Beam Eye Tracker die Trackingdaten mit anderen Programmen. Diese {\ac{API}} ermöglicht die Nutzung der Trackingdaten in eigenen Programmen. 

Die Beam {\ac{API}} kann über die Programmiersprachen C++, C, C\# und Python angesprochen werden \cite{lit_code_beam_sdk}. Im Rahmen dieser Projektarbeit wird die {\ac{API}} mit Pyhton angesprochen. Verschiedene Klassen in der {\ac{API}} ermöglichen eine individuelle Einbindung der Eye Tracking Daten in eine selbst entwickelte Anwendung. Relevant für diese Projektarbeit ist die Klasse \glqq{}screen\_gaze\_info\grqq{}. Diese beinhaltet die x- und y- Koordinaten und die Angaben, ob die Augen erkannt werden und wie gut die Schätzung ist.

Durch die anwenderfreundlichen {\ac{API}} des Beam {\ac{SDK}} und dem niedrigen Preis von 29,99€ des Eye Trackers fiel die Entscheidung, diesen für das Projekt zu nutzen. In Rücksprache mit den Beam Entwicklern wurde der Beam Eye Tracker für dieses Projekt unentgeltlich zur Verfügung gestellt.
>>>>>>> fd0befd8e95e2bae230d21ece8a1862eada6b577
